% =============================================================================
% Capítulo 3 — Desenvolvimento
% Começar o capítulo numa nova página e incluir uma breve descrição sobre
% qual será o assunto do capítulo.
% Este é o capítulo que descreve o desenvolvimento do sistema que constitui
% a solução implementada para o problema identificado.
% =============================================================================

\chapter{Desenvolvimento de \ldots}
\label{cap:desenvolvimento}

Este capítulo descreve o processo de desenvolvimento do trabalho, incluindo a análise de requisitos, a arquitetura do sistema e a implementação dos módulos desenvolvidos.

% ---- 3.1 Análise de Requisitos ----
\section{Análise de Requisitos}
\label{sec:requisitos}

% Quais são os processos, os dados a serem gerados pelo sistema, as restrições
% operacionais e as pessoas que irão utilizá-lo? De que forma esses elementos
% vão relacionar-se entre si? Os requisitos são as premissas para que o sistema
% a desenvolver possa cumprir os objetivos definidos para o sistema.

\subsection{Requisitos Funcionais}
\label{subsec:req_funcionais}

Listar os requisitos funcionais do sistema.

\begin{enumerate}
    \item RF01 --- O sistema deve permitir \ldots
    \item RF02 --- O sistema deve permitir \ldots
    \item RF03 --- O sistema deve permitir \ldots
\end{enumerate}

\subsection{Requisitos Não Funcionais}
\label{subsec:req_nao_funcionais}

Listar os requisitos não funcionais do sistema.

\begin{enumerate}
    \item RNF01 --- O sistema deve ser acessível \ldots
    \item RNF02 --- O sistema deve responder em menos de \ldots
\end{enumerate}

% ---- 3.2 Arquitetura do Sistema ----
\section{Arquitetura do Sistema}
\label{sec:arquitetura}

% Qual é a arquitetura definida para o sistema desenvolvido, com base nos
% requisitos? Usar um diagrama que identifique todos os módulos do sistema
% e os fluxos de dados entre esses módulos. Apresentar o diagrama seguido
% de uma descrição genérica.

Descrever a arquitetura do sistema, incluindo diagramas relevantes.

% Exemplo de inclusão de figura:
% \begin{figure}[htbp]
%     \centering
%     \includegraphics[width=0.8\textwidth]{arquitetura.png}
%     \caption{Arquitetura do sistema.}
%     \label{fig:arquitetura}
% \end{figure}

% ---- 3.3 Implementação do Módulo A ----
\section{Implementação do Módulo A}
\label{sec:modulo_a}

% Como é que o módulo A foi implementado? A descrição deve ser suficientemente
% detalhada para explicar todos os aspetos da implementação do módulo.
% Use diagramas, figuras, tabelas, equações, algoritmos, etc., quanto possível.

Descrever a implementação do módulo A.

% Heading 3 — \subsection
\subsection{Subsecção do Módulo A}
\label{subsec:modulo_a_sub}

Descrever os aspetos relevantes desta subsecção.

% Heading 4 — \subsubsection
\subsubsection{Sub-subsecção}
\label{subsubsec:modulo_a_subsub}

Descrever conteúdo de nível 4.

% Heading 5 — \paragraph (usado apenas quando estritamente necessário)
\paragraph{Parágrafo com título}

Descrever conteúdo de nível 5. Usar apenas quando absolutamente necessário.

% Exemplo de código fonte:
% \begin{lstlisting}[language=Python, caption={Exemplo de código.}, label={lst:exemplo}]
% def hello_world():
%     print("Hello, World!")
% \end{lstlisting}

% ---- 3.4 Implementação do Módulo B ----
\section{Implementação do Módulo B}
\label{sec:modulo_b}

% Como é que o módulo B foi implementado? Repetir a estrutura anterior
% para cada módulo identificado no diagrama de arquitetura do sistema.

Descrever a implementação do módulo B.

% ---- 3.5 Implementação do Módulo ... ----
% \section{Implementação do Módulo \ldots}
% \label{sec:modulo_x}
% Repetir para cada módulo adicional.

% ---- 3.6 Protótipo Final ----
\section{Protótipo Final}
\label{sec:prototipo}

% Se o projeto consistiu no desenvolvimento de algum tipo de protótipo
% (aplicativo de software ou dispositivo físico), apresentar e descrever
% esse protótipo final nesta secção.

Apresentar o protótipo ou as interfaces desenvolvidas.
